\justifying
\NSFCBodyText

\subsubsection{年度研究计划}

按照本次 CCF--百度松果基金研究资料，以\textbf{【暂定：合同生效后12个月】}组织任务，月份均为相对月份。经费与最终验收要求依据\textbf{【待补 F-PACT-01：正式课题附件及合同条款】}核定。以下安排体现各路线的依赖关系，不以模型规模或论文数量代替研究进展。

\textbf{第1--2个月：建立可比较的实验起点。}核验公开数据与模型许可，固定工具沙箱、任务序列和四类数据分区；复现静态低秩适配、基本验证器、经验回放和温度缩放。完成小规模成本测量，预先确定主要比较、覆盖率、遗忘容限和试验规模。形成任务协议、基础结果及成本清单，作为后续算法比较的共同起点。

\textbf{第3--4个月：检验任务证据与参数更新需求。}研究任务状态表征和有限探测更新，比较证据驱动分配与 AdaLoRA 等方法。在预算匹配后判断是否存在稳定增益；若任务证据未提供增量信息，则以固定配置保留可运行主线，并记录其边界。交付参数适配模块和相应消融结果。

\textbf{第5--6个月：区分轨迹正确性与训练价值。}建立结果与步骤证据，审计验证器错误及样本覆盖，完成同任务轨迹选择与偏好构造。对相关错误、反馈延迟及困难样本进行独立干预，形成可信轨迹学习模块和经版本标识的候选池。

\textbf{第7--8个月：研究有界记忆下的能力保持。}在固定适配器与轨迹记忆容量下比较复用、合并、重放及干扰约束，检验轨迹选择与历史记忆的相互作用。交付连续任务评价与遗忘分析；若合并导致历史能力明显下降，保留容量内的冻结策略并缩减扩展任务。

\textbf{第9--10个月：完成行动校准与双层验证。}对冻结策略建立独立校准器，比较固定、周期更新与漂移触发更新；开展影子审计、风险覆盖率分析及公开时序补充实验。完成核心流程在飞桨生态中的兼容检查，记录实际支持范围。

\textbf{第11--12个月：形成可复核结论与交付。}完成多种任务顺序、关键模型规模、消融与成本复核，整理失败案例和统计区间。冻结代码、配置、适配器与可公开轨迹，形成研究论文稿、技术报告和实验复现说明，按合同要求整理验收材料。

\subsubsection{预期研究结果}

预期形成一套 PACT-Agent 研究原型，包含参数适配、轨迹学习、有界记忆和行动校准四类模块；形成一套版本化任务及评价资源，完整记录训练、生成、验证、存储和推理成本；形成关于证据条件、任务变化与适应--遗忘--风险关系的实验分析。每项主要结论附对照、重复试验和失败条件，未观察到显著增益时也提交可解释的否定性结果。

拟形成可投稿的研究论文、开源代码及技术报告，具体数量与开放范围按任务书确定。可公开资产仅包含许可允许的代码、配置与轨迹；第三方模型权重、受限制数据及既有合同成果按原许可处理。预期成果不列未经依据的专利、性能百分比或真实业务收益承诺。
